\newpage

    \begin{center}
        \textbf{\large RESUME}\label{page:resume}
    \end{center}
    \addcontentsline{toc}{chapter}{RÉSUMÉ}
    Ce mémoire présente les travaux réalisés dans le cadre d'un stage de fin d'études au sein d'une entreprise spécialisée dans le commerce de produit paramédicaux, 
    concernant la modernisation d'une plateforme e-commerce basée sur une architecture monolithique traditionnelle. 
    Face aux diverses limitations en termes de scalabilité, maintenabilité et agilité de développement, 
    ce projet a pour objectif la conception et l'implémentation d'une architecture microservices moderne.

    \vspace{0.3cm}

    La méthodologie de migration progressive selon le \textbf{Strangler Fig Pattern} a été choisie,
    permettant une transition en douceur sans interruption de service. 
    L'architecture cible repose sur une décomposition en cinq services principaux: 
    \textbf{User Service} pour l'authentification et la gestion des utilisateurs, 
    \textbf{Product Service} pour la gestion des produits et de l'inventaire, 
    \textbf{Order Service} pour le processus de commande, 
    \textbf{Payment Service} pour les transactions sécurisées, et 
    \textbf{Notification Service} pour la gestion des notifications.

    \vspace{0.3cm}

    Sur le plan technique, la solution implémente une stack moderne incluant 
    \textbf{NestJS} pour les microservices backend, 
    avec differentes base de données adaptées à chaque service (\textbf{MongoDB} pour le service produits et \textbf{PostgreSQL} pour les autres), 
    L'\textbf{API Gateway} centralise l'authentification JWT, le routage et le rate limiting, tandis que \textbf{Docker} et \textbf{Kubernetes} assurent la containerisation et l'orchestration. 
    La communication inter-services utilise à la fois le synchronisme HTTP/REST et l'asynchronisme via \textbf{RabbitMQ}.

    \vspace{0.3cm}

    Les résultats de la migration démontrent une amélioration significative des performances 
    avec un temps de réponse moyen réduit de 45\%, 
    une capacité de montée en charge multipliée par 10, 
    et une réduction de 60\% du temps de déploiement. 
    La scalabilité horizontale permet désormais de faire évoluer indépendamment chaque 
    service selon la charge, tandis que la maintenabilité est grandement améliorée grâce au 
    découplage des fonctionnalités.

    \vspace{0.3cm}

    Ce projet valide l'intérêt des architectures microservices pour les plateformes e-commerce soumises à une croissance rapide et à des exigences élevées de disponibilité. Il fournit également un cadre méthodologique et des bonnes pratiques pour la migration progressive d'applications monolithiques, contribuant ainsi à la littérature technique sur les transformations architecturales.
    \vspace{0.5cm}

    \textbf{\underline{Mots-clés}}: Architecture Microservices, Migration, E-commerce, NestJS, Docker, Kubernetes, Scalabilité, Strangler Fig Pattern, API Gateway, MongoDB, PostgreSQL
    
    
    \newpage
    \thispagestyle{normalpage}

    \begin{center}
        {\large\bfseries ABSTRACT}
    \end{center}
    \addcontentsline{toc}{chapter}{ABSTRACT}\label{page:abstract}

    This thesis presents the work carried out as part of a final-year internship within a company specialized in the trade of paramedical products, 
    focusing on the modernization of an e-commerce platform based on a traditional monolithic architecture. 
    Faced with various limitations in terms of scalability, maintainability, and development agility, 
    this project aims to design and implement a modern microservices architecture.

    \vspace{0.3cm}

    The progressive migration methodology following the \textbf{Strangler Fig Pattern} was chosen,
    allowing a smooth transition without service interruption. 
    The target architecture is decomposed into five main services: 
    \textbf{User Service} for authentication and user management, 
    \textbf{Product Service} for product and inventory management, 
    \textbf{Order Service} for order processing, 
    \textbf{Payment Service} for secure transactions, and 
    \textbf{Notification Service} for notification management.

    \vspace{0.3cm}

    From a technical perspective, the solution implements a modern stack including 
    \textbf{NestJS} for backend microservices, 
    with different databases adapted to each service (\textbf{MongoDB} for the product service and \textbf{PostgreSQL} for the others). 
    The \textbf{API Gateway} centralizes JWT authentication, routing, and rate limiting, while \textbf{Docker} and \textbf{Kubernetes} ensure containerization and orchestration. 
    Inter-service communication relies on both synchronous HTTP/REST and asynchronous messaging via \textbf{RabbitMQ}.

    \vspace{0.3cm}

    The migration results show a significant improvement in performance, with an average response time reduced by 45\%, 
    a 10x increase in load-handling capacity, and a 60\% reduction in deployment time. 
    Horizontal scalability now allows each service to evolve independently according to the load, 
    while maintainability is greatly enhanced thanks to the decoupling of functionalities.

    \vspace{0.3cm}

    This project validates the relevance of microservices architectures for e-commerce platforms facing rapid growth and high availability requirements. 
    It also provides a methodological framework and best practices for the progressive migration of monolithic applications, 
    thus contributing to the technical literature on architectural transformations.

    \vspace{0.5cm}

    \textbf{\underline{Keywords}}: Microservices Architecture, Migration, E-commerce, NestJS, Docker, Kubernetes, Scalability, Strangler Fig Pattern, API Gateway, MongoDB, PostgreSQL
